% BANKS-SOURCE: pdf=109 print=99 kind=section
\subsection{A brief history of the physics of non-abelian gauge theory}

% BANKS-SOURCE: pdf=109 print=99 kind=prose
A resource for the history of the standard model is \cite{banks-ref-55}. Non-abelian gauge
theory is an integral part of Einstein's theory of gravitation, as was realized
by Cartan \cite{banks-ref-56}. Weyl \cite{banks-ref-57} introduced the term gauge symmetry in an attempt to
explain electromagnetism in terms of the scale ambiguity of the space-time
metric (Weyl transformations). The first inkling that it had something to do
with particle physics came in the Fermi theory of weak interactions.\footnote[4]{Curiously,
in 1938 Oskar Klein wrote down a Lagrangian similar to the standard electroweak
theory \cite{banks-ref-58,banks-ref-59}, but did not use his knowledge of Kaluza--Klein compactification
to make it gauge-invariant.} Fermi postulated the existence of a charge-carrying
vector

% BANKS-SOURCE: pdf=110 print=100 kind=figure id=8.1
\begin{figure}[htbp]
  \centering
  % BANKS-SOURCE: pdf=110 print=100 kind=figure id=8.1
\begin{tikzpicture}[x=1cm,y=1cm,
    BanksFermion/.style={line width=.8pt,postaction={decorate},
      decoration={markings,mark=at position #1 with
        {\arrow{Stealth[length=4.6pt,width=4.4pt]}}}},
    BanksFermion/.default=.5,
    BanksW/.style={line width=.8pt,decorate,
      decoration={snake,amplitude=6.2pt,segment length=9pt}}]
  \coordinate (v1) at (0,0);
  \coordinate (v2) at (1.3,0);
  % p comes in from the upper left, n leaves towards the lower left
  \draw[BanksFermion=.52] (-1.07,1.73) -- (v1);
  \draw[BanksFermion=.52] (v1) -- (-1.09,-1.72);
  % antineutrino comes in from above, electron leaves to the upper right
  \draw[BanksFermion=.47] (1.10,1.73) -- (v2);
  \draw[BanksFermion=.50] (v2) -- (1.95,1.70);
  % the W boson, exchanged between the two vertices
  \draw[BanksW] (v1) -- (v2);
  \fill (v1) circle (2.3pt);
  \fill (v2) circle (2.3pt);
  \node[above] at (-1.07,1.73) {$p$};
  \node[below] at (-1.09,-1.72) {$n$};
  \node[above] at (1.10,1.73) {$\bar\nu_e$};
  \node[above] at (1.95,1.70) {$e$};
  \node[below] at (0.81,-0.24) {$W^+$};
\end{tikzpicture}

  \caption{Fermi's theory of weak interactions.}
  \label{fig:8.1}
\end{figure}

% BANKS-SOURCE: pdf=110 print=100 kind=prose
boson $W_\mu^+$, which could mediate the process of neutron decay via the
Feynman diagram of Figure~\ref{fig:8.1}.

Fermi, the father of the idea of effective field theory, realized that, at
distances long compared with the Compton wavelength of the $W$, this diagram
would be the same as an effective four-fermion interaction:
The Fermi interaction in this historical discussion is four-dimensional.
\begin{equation*}
  \mathcal L_{\mathrm{eff}}=G_F(\bar p\gamma^\mu n\,\bar\nu\gamma_\mu e),
\end{equation*}
and this is what came to be known as the Fermi theory of weak interactions.
Nuclear data were soon shown to be inconsistent with this simple form, and
$\gamma^\mu$ was replaced with an arbitrary sum of Dirac matrices. It was not
until the mid 1950s, after the ground-breaking work of Lee and Yang \cite{banks-ref-60} had
showed that weak interactions did not preserve parity, that Marshak and
Sudarshan, Gershtein and Zel'dovich, and Feynman and Gell-Mann \cite{banks-ref-61,banks-ref-62,banks-ref-63}
realized that, if one allowed the $W$ boson to couple to a linear combination
of vector and axial currents, then one could construct an elegant theory, which
also fit the data. The stage was set for a non-abelian gauge theory of weak
interactions.

% BANKS-SOURCE: pdf=110 print=100 kind=prose
In 1954, Yang and Mills \cite{banks-ref-64} had introduced non-abelian gauge theory into
particle physics in an attempt to account for the strongly interacting vector
bosons.\footnote[5]{The modern approach to this idea is outlined in Problem
8.15.} Bludman \cite{banks-ref-65} and Schwinger \cite{banks-ref-66} realized that one could adapt this
technology to the Fermi theory and that an $SU(2)$ gauge theory could unify the
weak and electromagnetic interactions. There were two problems with this idea.
Glashow \cite{banks-ref-67} showed that the weak leptonic currents did not close on an algebra
with the electromagnetic current (Problem 8.14). One either had to introduce
new leptons or follow the route chosen by Glashow (and the world) of introducing
an $SU(2)\times U(1)$ gauge group, with the photon associated with a linear
combination of one of the $SU(2)$ generators and the $U(1)$.

The remaining problem, that of the mass of the non-abelian bosons, was not so
easily solved. Mass terms seemed to break the non-abelian gauge symmetry.
Furthermore, in four-dimensional Euclidean space, the propagator of a massive
vector boson behaves like
\begin{equation*}
  \frac{\delta_{\mu\nu}+p_\mu p_\nu/m^2}{p^2+m^2},
\end{equation*}
and the effect of the mass does not appear to vanish at high momentum.

% BANKS-SOURCE: pdf=111 print=101 kind=prose
In the meantime, developments stemming from the theory of superconductivity
were leading to a solution of the problem. Following work of Schwinger on
two-dimensional quantum electrodynamics \cite{banks-ref-68}, Anderson \cite{banks-ref-69} discussed the
screening of electromagnetic interactions in condensed matter systems. Higgs
showed how to incorporate the Meissner effect into a fully relativistic
treatment of electrodynamics \cite{banks-ref-70,banks-ref-71,banks-ref-72}. We will study his model in the next
section. Developments followed rapidly \cite{banks-ref-73,banks-ref-74}. Weinberg \cite{banks-ref-75} and Salam \cite{banks-ref-76}
combined the Higgs mechanism with Glashow's $SU(2)\times U(1)$ theory to
construct the modern theory of the electro-weak interaction. A few years later,
't Hooft and Veltman \cite{banks-ref-77,banks-ref-78} proved that spontaneously broken gauge theories,
as the new theories came to be called, were renormalizable.

I was a graduate student at the time of these developments, and it was obvious
that a generalization of these ideas to the strong interactions was called for.
In fact, such a generalization already existed, though it was relatively
obscure, particularly in Europe. It was not the original Yang--Mills theory of
vector bosons, but rather Nambu's $SU(3)$ color gauge theory \cite{banks-ref-79} invented to
explain the statistics of quarks. In the quark model, baryons are bound states
of three quarks, and the correct spectrum of baryons is obtained if one chooses
a wave function symmetric under interchange of the three quarks. This is
paradoxical because the quarks carry spin $1/2$. Greenberg \cite{banks-ref-80} realized that
the paradox could be removed if one introduced an additional three-valued label
for the quarks (color) and insisted that all states be singlets under an $SU(3)$
group transforming this label.\footnote[6]{Greenberg actually used the language
of parastatistics rather than color. The two ideas are equivalent for classifying
possible states but dynamically different. In modern language, parastatistics is
the impossible limit in which the QCD scale is taken to infinity with hadron
masses kept fixed.} Nambu \cite{banks-ref-79} invented the color gauge theory as an attempt to
explain the color singlet condition dynamically. He showed that the
single-gauge-boson exchange forces between quarks would lower the energy of the
singlet states (just as neutral states have less electromagnetic energy than
charged states). In Nambu's picture, the colored states, including quarks, would
eventually be found at higher energies. It was perhaps for this reason that he
also figured out a way to give the quarks integer electric charge \cite{banks-ref-81}, using
the new degree of freedom. In the modern view of quantum chromodynamics or QCD,
Nambu's calculation is just giving the short-distance part of the interquark
forces. At long distances, the potential rises linearly and quarks are
permanently confined. Such a potential cannot arise from particle exchange. To
understand it, we will turn in the next section to the Higgs theory of
relativistic superconductivity, which provides both the framework for
electro-weak physics and (after a duality transformation) the explanation of
quark confinement.
